\chapter{The Woman on Platform 8}
I was my second year at boarding school, and I was sitting on platform no. 8 at Ambala station, waiting for the northern bound train. I think I was about twelve at the time. My parents considered me old enough to travel alone, and I had arrived at Ambala by bus early in the evening, now there was a wait till midnight before my train arrived. Most of the time, I had been pacing up and down the platform, browsing at the bookstall, or feeding broken biscuits to stray dogs; trains came and went, and the platform would be quiet for a while and then, when a train arrived, it would be an inferno of heaving, shouting, agitated human bodies. As the carriage doors opened, a tide of people would swoop down upon the nervous little ticket collector at the gate; and every time this happened, I would be caught in the rush and swept outside the station. Now tried of this game and of ambling about the platform, I sat down on my suitcase and gazed dismally across the railroad tracks.
\par
Trolleys rolled past me, and I was conscious of the cries of the various vendors-the men who sold curds and lemon, the sweetmeat seller, the newspaper boy-but I had lost interest in all that went on along busy platform, and continued to stare across the railroad tracks, feeling bored and a little lonely.
\par
'Are you all alone, my son?' asked a soft voice close behind me.
\par
I looked up and saw a woman standing near me. She was learning over and I saw a pale face and dark, kind eyes. She wore no jewels, and was dressed very simply in a white sari. She seemed poor.
\par
'Yes, I am going to school,' I said, and stood up respectfully.
\par
'I have been watching you for some time,' she said. 'Didn't you parents come to see you off?'
\par
'I don't live here' I said. 'I had to change trains. Anyway, I can travel alone.'
\par
'I am sure you can,' she said, and I liked her for saying that, and I also liked her for the simplicity of her dress, and for her deep, soft voice and the serenity of her face.
\par
Tell me, what is your name?' she asked.
\par
'Arun,' I said.
\par
'And how long do you have to wait for your train?'
\par
'About an hour, I think. It comes at twelve o' clock.
\par
'Then come with me and have something to eat.'
\par
I was going to refuse, our of shyness and suspicion, but she took me by the hand, and then I felt it would be silly to pull my hand away. She told a coolie to look after my suitcase, and then she led me away down the platform. Her hand was gentle, and she held mine neither too firmly not too lightly. I looked up at her again. She was not young. And she was not old. She must have been over thirty but, have she been fifty, I think she would have looked much the same.
\par
She took me into the station dining room, ordered tea and samosas and jalebis, and at once I began to thaw and take a new interest in this kind woman. The strange encounter had little effect on my appetite. I was hungry schoolboy, and I ate as much as I could in as polite a manner as possible. She took obvious pleasure in watching me eat, and I think it was the food that strengthened the bond between us and cemented our friendship, for under the influence of the tea and sweets, I began to talk quite freely, and told her about my school, my friends, my likes and dislikes. She questioned me quitely form time to time, but seemed to prefer listening; she drew me out very well, and I had soon forgotten that we were strangers. But she did not ask her where she lived. I accepted her for what she had been to me-a quiet, kind and gentle woman who gave sweets to a lonely boy on a railway platform . . .
\par
After about half an hour, we left the dining room and began walking back along the platform. An engine was shunting up and down beside platform no. 8, and as it approached, a boy leapt off the platform and ran across the rails, taking a shortcut to the next platform. He was at a safe distance from the engine, and there was no danger unless he had fallen; but as he leapt across the rails, the woman clutched my arm. Her fingers due into my flesh and I winced with pain. I caught her fingers and looked up at her; I saw a spasm of pain fear and sadness pass across her face. She watched the boy as he climbed the other platform, and it was not until he had disappeared in the crowd that she relaxed her hold on my arm. She smiled at me reassuringly and took my hand again, but her fingers trembled against mine.
\par
'He was all right,' I said, feeling that if was she who needed reassurance.
\par
She smiled gratefully at me and pressed my hand. We walked together in silence until we reached the place where I had left my suitcase, one of my schoolfellows, Satish, a boy of about my age, had turned up with his mother.
\par
'Hello, Arun!' he called. 'The train's coming in late as usual. Did you know we have a new headmaster this year?'
\par
We shook hands, and then he turned to his mother and said, 'This is Arun, Mother. He is one of my friends and the best bowerler in the class.'
\par
'I am glad to know that,' said his mother, a large imposing woman who wore spectacles. She looked at the woman who led my hand and said, 'And I suppose you're Arun's mother?'
\par
I opened my mouth to make some explanation, but before I could say anything the woman replied, 'Yes, I am Arun's mother.'
\par
I was unable to speak a word. I looked quickly up at the woman, but she did not appear to be at all embarrassed and was smiling at Satish's mother.
\par
Satish's mother said, 'It's such a nuisance having to wait for the train right in the middle of the night. But one can't let the child wait here station like this, there are so many suspicious characters hanging about. These days one has to be very careful fo strangers.'
\par
'Arun can travel alone though,' said the woman beside me, and somehow I felt grateful to her for saying that. I had already forgiven her for lying and besides, I had taken an instinctive dislike to Satish's mother.
\par
'Well, by very careful, Arun,' said Satish's mother, looking sternly at me through her spectacles. 'Be very careful when your mother is not with you. And never talk to strangers!'
\par
I looked from Satish's mother to the woman who had given me tea and sweets, and then back at Satish's mother.
\par
'I like strangers,' I said.
\par
Satish's mother definitely staggered a little, as obviously she was not used to being contradicted by small boys. 'There you are, you see! If you don't watch over them all the time, they'll walk straight into trouble. Always listen to what your mother tells you,' she said, wagging a fat little finger at me. 'And never, never talk to strangers.'
\par
I glared resentfully at her and moved closer to the woman who had befriended me. Satish was standing behind his mother, grinning at me, and delighting in my clash with his mother. Apparently, he was on my side.
\par
The station bell clanged, and the people who had till now been squatting resignedly on the platform began bustling about.
\par
'Here it comes,' shouted Satish, as the engine whistle shrieked and the front lights played over the rails.
\par
The train moved slowly into the station, the engine hissing and sending out waves of stream. As it came to stop, Satish jumped on the footboard of a lighted compartment and shouted, 'Come on, Arun, this one's empty!' and I picked up my suitcase and made a dash for the open door. 
\par
We placed ourselves at the open windows, and the two women stood outside on the platform, talking to us. Satish's mother did most of the talking.
\par
'Now don't jump on and off moving trains, as you did just now,' she said. 'And don't stick your heads out of the windows and don't eat any rubbish on he way.' She allowed me to share the benefit of her advice, as she probably didn't think my 'mother' a very capable person. She handed Satish a bag of fruit, a cricket bat and a big box of chocolates, and told him to share the food with me. Then she stood back from the window to watch how my 'mother' behaved.
\par
I was smarting under the patronizing tone of Satish's mother, who obviously thought mine a very poor family, and I did not intend to give the others woman away. I let her take my hand in hers, but I could think of nothing to say. I was conscious of Satish's mother staring at us with hard, beady eyes, and I found myself hating her with a firm, unreasoning hate. The gerund walked up the platform, blowing his whistle for the train to leave. I looked straight into the eyes of the woman who held my hand, and she smiled in a gentle, understanding way. I learned out of the window then, and put my lips to her cheek and kissed her.
\par
The carriage jolted forward, and she drew her hand away.
\par
'Goodbye,' I said to the other woman, 'Goodbye . . . mother . . .'
\par
I didn't wave or shout, but sat still in the front of the window, gazing at the woman on the platform. Satish's mother was talking to her, but she didn't appear to be listening; she was looking at me as the train took me away. She stood there on the busy platform, a pale sweet woman in white, and I watched her until she was lost in the milling crowd.
